% subsubsection: 極形式による図形的解釈

\subsubsection{極形式による図形的解釈}
  \BookFigMobius

  図では,点$a_n$と二つの\bookterm{fixed}{不動点}$\alpha,\beta$を結ぶ線分を描いた.
  上の線分の長さを下の線分の長さで割ると,
  \[
    \frac{|a_n-\alpha|}{|a_n-\beta|}
    =\left|\frac{a_n-\alpha}{a_n-\beta}\right|
    =|b_n|
  \]
  となる.また,二つの線分の向きの差は,
  $b_n$の\bookterm{argument}{偏角}$\arg b_n$に対応する.
  したがって,図で見ている\textbf{長さの比と角度の差}は,
  計算で導入する一つの複素数$b_n$の絶対値と\bookterm{argument}{偏角}に他ならない.
  さらに$b_{n+1}=\lambda b_n$となったときは,
  長さの比が$|\lambda|$倍され,角度の差が$\arg\lambda$だけ増える.
  特に$|\lambda|=1$なら,長さの比を保ったまま\bookterm{argument}{偏角}だけが変わり,
  \bookterm{complex-plane}{複素平面}上の回転として理解できる.


  ここまででは,\,$\alpha,\beta$を\bookterm{fixed}{不動点}として
  \[
    b_n=\frac{a_n-\alpha}{a_n-\beta}
  \]
  と置くことで,1次分数型の\bookterm{recurrence}{漸化式}を\bookterm{geometric}{等比数列}に帰着できることを示した.

  ここでは,この変形を\bookterm{complex-plane}{複素平面}上で考えてみよう.
  まず注意したいのは,
  \[
    \frac{a_n-\alpha}{a_n-\beta}
  \]
  そのものを「距離の比」と呼ぶのは正確ではないという点である.
  実数上でも符号の情報を含み,\bookterm{complex-plane}{複素平面}上では\bookterm{argument}{偏角}の情報まで含む.
  実際に距離の比を表すのは絶対値を取った
  \[
    \left|\frac{a_n-\alpha}{a_n-\beta}\right|
    =
    \frac{|a_n-\alpha|}{|a_n-\beta|}
  \]
  であり,これは点$a_n$から$\alpha,\beta$までの距離の比である.
  一方,もとの複素数の比には,この距離の比に加えて二つの差の\bookterm{argument}{偏角}の差も記録されている.

  特に,\,$\alpha,\beta$が実数ではなく複素数となる場合には,この変形によって\bookterm{recurrence}{漸化式}を「回転」として解釈することができる.

  $p,q,r,s$が実数であるとき,
  \[
    rx^2+(s-p)x-q=0
  \]
  の2解が実数でなければ,解と係数の関係から,2解は互いに\bookterm{conjugate}{共役}な複素数となる.
  そこで,その2解を
  \[
    \beta=\overline{\alpha}
  \]
  とする.

  このとき,
  \[
    b_n=\frac{a_n-\alpha}{a_n-\overline{\alpha}}
  \]
  と置くと,先ほどと同じ変形によって
  \[
    b_{n+1}
    =
    \frac{p-r\alpha}{p-r\overline{\alpha}}b_n
  \]
  となる.

  ここで,
  \[
    \lambda=\frac{p-r\alpha}{p-r\overline{\alpha}}
  \]
  とおく.
  \, $p,r$は実数であり,\,$\overline{p-r\alpha}=p-r\overline{\alpha}$であるから,
  \[
    |\lambda|
    =
    \left|
    \frac{p-r\alpha}{p-r\overline{\alpha}}
    \right|
    =1
  \]
  である.

  したがって,\,$\lambda$は\bookterm{polar}{極形式}を用いて
  \[
    \lambda=\cos\theta+i\sin\theta
  \]
  と表すことができる.

  一方,\,$b_n$も\bookterm{polar}{極形式}を用いて
  \[
    b_n=\rho_n(\cos\varphi_n+i\sin\varphi_n)
  \]
  と表すことができる.
  このとき,
  \begin{align*}
    b_{n+1}
    &=
    (\cos\theta+i\sin\theta)
    \rho_n(\cos\varphi_n+i\sin\varphi_n)\\
    &=
    \rho_n
    \{\cos(\varphi_n+\theta)
    +i\sin(\varphi_n+\theta)\}
  \end{align*}
  であるから,
  \[
    \rho_{n+1}=\rho_n,
    \qquad
    \varphi_{n+1}=\varphi_n+\theta
  \]
  となる.

  つまり,$b_n$は\bookterm{complex-plane}{複素平面}上で原点からの距離を変えずに,\,$\theta$だけずつ回転している.

  このことから,
  \[
    b_{n+1}=\lambda b_n
  \]
  という\bookterm{geometric}{等比数列}への帰着は,単なる計算上の技巧ではなく,\bookterm{complex-plane}{複素平面}上の図形的な変化を表していることが分かる.

  実際に,次の\bookterm{recurrence}{漸化式}を考えてみよう.
  \[
    a_{n+1}=\frac{a_n+1}{1-a_n}
  \]
  この\bookterm{recurrence}{漸化式}の\bookterm{fixed}{不動点}は
  \[
    x=\frac{x+1}{1-x}
  \]
  を解けばよいから,
  \[
    x^2+1=0
  \]
  より
  \[
    \alpha=i,\qquad\beta=-i
  \]
  である.

  そこで,
  \[
    b_n=\frac{a_n-i}{a_n+i}
  \]
  と置くと,
  \begin{align*}
    b_{n+1}
    &=
    \frac{1+i}{1-i}b_n
  \end{align*}
  となる.

  ここで,
  \[
    \frac{1+i}{1-i}
    =
    \frac{(1+i)^2}{2}
    =i
  \]
  であるから,
  \[
    b_{n+1}=ib_n
  \]
  を得る.

  $i$は\bookterm{polar}{極形式}で
  \[
    i=\cos\frac{\pi}{2}+i\sin\frac{\pi}{2}
  \]
  と表せる.
  したがって,\,$b_n$は\bookterm{complex-plane}{複素平面}上で毎回反時計回りに$\dfrac{\pi}{2}$だけ回転する.

  例えば,\,$a_1=\dfrac12$とすると,
  \begin{align*}
    a_1&=\frac12\\
    a_2&=3\\
    a_3&=-2\\
    a_4&=-\frac13\\
    a_5&=\frac12
  \end{align*}
  となり,\,$a_n$は
  \[
    \frac12\longrightarrow3\longrightarrow-2
    \longrightarrow-\frac13\longrightarrow\frac12\longrightarrow\cdots
  \]
  と変化する.

  一方,\,$b_n$に着目すると,
  \[
    b_{n+1}=ib_n
  \]
  であるから,
  \[
    b_1\longrightarrow ib_1\longrightarrow-b_1
    \longrightarrow-ib_1\longrightarrow b_1\longrightarrow\cdots
  \]
  となる.

  つまり,実数上では複雑に見える
  \[
    a_1,a_2,a_3,\ldots
  \]
  の変化が,
  \[
    b_n=\frac{a_n-i}{a_n+i}
  \]
  という変換によって,\bookterm{complex-plane}{複素平面}上で一定角度ずつ回転する点の運動として表される.

  このように,異なる二つの\bookterm{fixed}{不動点}をもつ1次分数型\bookterm{recurrence}{漸化式}では,
  \[
    \frac{a_n-\alpha}{a_n-\beta}
  \]
  という変換を用いることで,\bookterm{recurrence}{漸化式}を\bookterm{geometric}{等比数列}に変形できる.
  特に,係数が実数で\bookterm{fixed}{不動点}が非実の\bookterm{conjugate}{共役}な二点であるときには,変換後の点の変化を\bookterm{complex-plane}{複素平面}上の回転として図形的に解釈することができる.
